\scanpage{110}
\mspage{25}
% archive sheet number: 109
% Lines 1--6.
La première assertion résulte \inserted{aussitôt} de la détermination du ``groupe fondamental modéré'' de $U$, v. plus bas (\quad). Pour la deuxième, on procède par récurrence sur $r$, le cas $r=0$ étant trivial, et le cas $r=1$ \deleted{\unsure{est}} \unsure{étant} est contenu dans 4.2 (où on fait $d=1$). \deleted{\unsure{Si}} Si $r\geq2$, on \deleted{écrit} considère

\[
\begin{aligned}
U'&=X-\deleted{\unsure{\sum|D_i}}\bigcup_{1\leq i\leq r-1}D_i\\
D'_r&=\deleted{\unsure{X\cap D_r}}U'\cap D_r
\end{aligned}
\]

de sorte que

\[
\begin{aligned}
U&=U'-D'_r\ ;\ \deleted{\unsure{D'_r=D_r-\sum_{1\leq i\leq r-1}D_i\cap D_r}}\\
D'_r&=D_r-\operatorname{supp}\Delta_r,\qquad \Delta_r=\sum_{1\leq i\leq r-1}D_i\cap D_r.
\end{aligned}
\]

\margin{\deleted{\unsure{la fibrée|la fibre}}}

On a la suite \unsure{ex.act|ex.acte} de cohomologie relative (4.3.1) pour $D'_r$ dans $U'$

% From this display to the bottom, the entire block is diagonally cancelled.
\[
\deleted{
\cdots\longrightarrow H^{i-2}_{\deleted{\unsure{D'_r}}}(D'_r,\deleted{\unsure{\Lambda}}\mu_n^{\unsure{\otimes-1|\otimes1}})\xrightarrow{\alpha_i}H^i(U',\Lambda)\longrightarrow H^i(U,\Lambda)
\longrightarrow H^{i-1}_{\deleted{D'_r}}(D'_r,\deleted{\unsure{\Lambda}}\mu_n^{\unsure{\otimes-1|\otimes1}})\longrightarrow\cdots.
}
\tag*{\deleted{(4.5.5)}}
\]

\[
\deleted{\inserted{H^{i-2}(D'_r,\Lambda)\otimes\mu_n(\unsure{k})^{-1}}}
\]

\deleted{D'autre part \deleted{\unsure{on a}} \unsure{la} suite spectrale}

\[
\deleted{H^{\unsure{*|\bullet}}_{D'_r}(U,\Lambda_U)\Longleftarrow \inserted{E_2^{pq}=}H^p(\unsure{D_r|D'_r},\underline{H}^{q}_{D'_r}(\deleted{\unsure{\mu_n}}\Lambda_{U'}))}
\]

\deleted{et la détermination 4.2 des $\underline{H}^q_{D'_r}(\Lambda_{U'})$ donne}

\[
\deleted{H^i_{D'_r}(U',\Lambda_{U'})\simeq H^{i-2}(D'_r,\mu_{nD'_r}^{\otimes-1})\simeq H^{i-2}(D'_r,\Lambda)\otimes\mu_n(k)^{-1}}
\]

\deleted{où $k$ est le corps résiduel de $X$.}

% A rewritten paragraph is inserted above the spectral-sequence paragraph; its continuation runs down the left edge.
\deleted{\inserted{L'hypothèse de récurrence, appliquée à $D'_r$ et les $r-1$ div. réguliers $D'_i$ $(1\leq i\leq r-1)$ sur $D'_r$}}

\deleted{\deleted{\unsure{D'autre part}, l'hypothèse de récurrence appliquée à $D'_r$ et aux \inserted{$r-1$ diviseurs réguliers} $(D_i\cap D_r)$ $(1\leq i\leq r-1)$ sur $D'_r$} implique:}

\[
\deleted{H^{\deleted{\unsure{*}}i}(D'_r,\Lambda)\simeq\bigwedge^{\deleted{\unsure{*}}i}H^1(D'_r,\Lambda)\simeq\left(\bigwedge^{\deleted{\unsure{*}}i}(\Lambda^{r-1})\right)\otimes\mu_n^{\otimes-i}(k)}
\]

\deleted{et de même l'hypothèse de récurrence appliquée à $X$ et aux $D_i$, $(1\leq i\leq r-1)$ donne}

\[
\deleted{H^i(U',\Lambda)\simeq\bigwedge^iH^1(U',\Lambda)\simeq\left(\bigwedge^i(\Lambda^{r-1})\right)\otimes\mu_n^{\otimes-i}(k).}
\]

\deleted{\deleted{Il faut exprimer \unsure{maintenant} l'hom de \unsure{(***)}}}

\[
\deleted{\deleted{H^i_{D'_r}(U,\Lambda)\simeq H^{i-2}(D'_r,\Lambda)\otimes\mu_n^{\unsure{\otimes1|\otimes-1}}(k)\simeq\bigwedge^{i-2}(\Lambda^{r-1})\otimes\mu_n^{\deleted{\unsure{\otimes(-i+1)|\otimes(-i)}}\inserted{\unsure{-(i-2)-1}}}(k)}}
\]

% A bent annotation arrow leads from the preceding first term to this added formula.
\[
\deleted{\inserted{H^i(U',\Lambda)\simeq\bigwedge^i(\Lambda^{r-1})\otimes\mu_n^{\otimes-i}(k)}}
\]
